\documentclass[11pt]{article}
\usepackage[a4paper,margin=22mm]{geometry}
\usepackage[T1]{fontenc}
\usepackage{lmodern,microtype,amsmath,amssymb,graphicx,caption,booktabs,array}
\usepackage[super,sort&compress]{natbib}
\usepackage{hyperref}
\hypersetup{colorlinks=true,urlcolor=blue,citecolor=blue,linkcolor=blue}
\widowpenalty=10000\clubpenalty=10000
\setlength{\parskip}{5pt}\setlength{\parindent}{0pt}\emergencystretch=4em
\title{Choosing scientific visualizations through task-first optimization}
\author{Mohamed Elgendi\\Khalifa University, Abu Dhabi, United Arab Emirates}
\date{}
\begin{document}\maketitle
\begin{abstract}
A visualization is useful for a question, not universally optimal for a dataset. We present an executable task-first framework that declares the requested information, identifies compatible chart implementations, recovers the target from rendered pixels and ranks task-specific errors. The evaluation spans 13 familiar chart families and six questions: median differences, spread differences, distributional differences, category shares, ordered profiles and association. The original benchmark contains 127 two-group comparisons from nine public resources; an extension adds 149 empirical task instances derived from existing resources and 48 synthetic compositions. Under one fixed setting per familiar family, box plots belong to the winning set in 19 of 127 median comparisons and 30 spread comparisons. Pie charts win 16 of 24 empirical composition cases and donuts the remaining eight under the primary observer. Eight worked examples demonstrate eight genuinely selected families, while a second set exposes changes after transferring frozen selections to a different display size. Results depend on the task, candidate budget and computational observer; specialized task-matched controls can outperform the familiar set. The contribution is a reproducible method for making those dependencies explicit and auditable, not a universal chart ranking or a validated model of human comprehension.
\end{abstract}

\section*{Introduction}
Scientists choose charts to answer questions: how groups differ, what share belongs to a category, how a measurement evolves or how two variables are related. A box plot, pie chart and scatter plot summarize different structures. Comparing them without first declaring the question can produce an authoritative-looking score for an ill-defined choice. A useful selector must explain what it is trying to preserve and which alternatives can meaningfully represent it.

Task taxonomies distinguish the purpose of a visualization from the means used to construct it.\cite{amar2005,brehmer2013} Automated graphical design, constraint systems and learned recommenders address complementary aspects of construction and recommendation.\cite{mackinlay1986,draco,vizml} Human graphical-perception studies evaluate judgments made by people.\cite{cleveland1984,heer2010} Our computational problem is narrower: given a numerical target and a declared pixel-recovery algorithm, which implemented chart best preserves that target under specified display conditions? Such a score can support an audit, but does not replace evidence about readers. Scientific visualization training should make this distinction visible.\cite{elgendi2017,elgendi2020}

An initial distributional benchmark illustrates the problem. Asking for five quantile differences favors encodings that preserve those quantities and can penalize a box plot for information it does not explicitly show. A pie chart has no natural role until a composition and a meaningful denominator have been defined. Restricting every comparison to one question would therefore misrepresent the scope of familiar charts. Conversely, adding inappropriate charts simply to reach a desired catalogue size would not make the comparison fair.

We address this with a task-to-chart contract: an input structure, target, candidate implementations, calibration, recovery rule and loss are declared before computing the selection. The contract is explicit enough to reproduce or challenge. The revised evaluation contains six questions and 13 familiar families, with one fixed setting per family for the cross-task audit. A separate 21-setting distributional search is retained to expose tuning-budget effects. We do not force each family to win or adjust scores to break genuine ties.

The figures follow this argument. Figure~\ref{fig:scoring} traces one score from a visible group difference to the arithmetic. Figure~\ref{fig:improvement} presents eight different selected families, including a box plot, pie chart, line chart and scatter plot. Figure~\ref{fig:limits} asks whether those kinds of selections survive a new display condition. Figure~\ref{fig:tasks} shows every familiar family evaluated and distinguishes a tested non-winner from an unimplemented task contract. The examples are illustrations; complete task-wise results determine the scope of the conclusions.

\section*{Results}
\subsection*{The question defines the optimization problem}
Let $D$ denote the input data and $t$ the question. The contract specifies a target $T_t(D)$, a candidate set $\mathcal C_t$, a rendering $I_{c,e,r}$ and public calibration $A_c$. A decoder $\mathcal D_{c,t}$ receives the image and calibration, but neither the source observations nor the target values. For environment $e$ and placement $r$, define
\begin{align}
\widehat T_{c,t,e,r}&=\mathcal D_{c,t}(I_{c,e,r},A_c),\\
\overline L_{c,t,e}&=\frac14\sum_{r=1}^{4}\ell_t\!\left(\widehat T_{c,t,e,r},T_t(D)\right),\\
S_{c,t,e}&=\frac{100}{1+\overline L_{c,t,e}},\qquad
\mathcal C_t^*(D,e)=\operatorname*{arg\,min}_{c\in\mathcal C_t}\overline L_{c,t,e}.
\label{eq:score}
\end{align}
Maximizing the score is exactly equivalent to minimizing the declared loss. The reciprocal transformation is an interpretation aid, not a new optimization theorem. Losses tied within $10^{-12}$ remain tied. Failed decodes invalidate that setting in the relevant environment; unimplemented task contracts receive no score, not zero.

Eligibility concerns the declared input and representation. A composition requires nonnegative categories with a meaningful total; an ordered profile requires an order; association requires paired variables. Rebinning a continuous variable into categories changes the question and is recorded as a transformation. Compatibility does not imply that an encoding contains all requested information: lossy summaries, including mean bars for quantile questions, remain explicitly labeled diagnostic controls. Their occasional numerical success does not prove general sufficiency.

For the marginal tasks, let $d_p=Q_y(p)-Q_x(p)$ and let $s$ be pooled within-group SD. Targets are $d_{.5}$ for the median, $d_{.75}-d_{.25}$ for interquartile-spread contrast, and $(d_{.1},d_{.25},d_{.5},d_{.75},d_{.9})$ for five distributional gaps. Their losses are absolute or mean absolute error divided by $s$. These are empirical distributional questions, not causal effects, diagnostic accuracy or within-person changes. Alternatives to mean-only summaries are particularly relevant when distributions differ beyond their means.\cite{weissgerber2015,rousselet2017}

Composition loss is total variation, $\frac12\sum_k|\widehat p_k-p_k|$, between recovered and true shares. Profile loss is mean absolute error across 12 successive bin means divided by their sample SD. Association loss is $|\widehat\rho-\rho|/2$, where $\rho$ is the Pearson correlation of the displayed paired subset. These losses answer different questions. A score of 99 for shares is not equivalent in difficulty or scientific importance to 99 for a median; scores and gains are compared within tasks only.

\subsection*{A worked example exposes both the data difference and the calculation}
Figure~\ref{fig:scoring} compares mean cell-nucleus area from 357 benign and 212 malignant WDBC records. The malignant group contains larger and more dispersed values. In the 21-setting familiar-chart search, raw dots score 97.70 versus 60.17 for mean bars. Average absolute recovery error is 5.85355 source units; dividing by pooled SD 248.3973 gives 0.023565, hence $100/(1+0.023565)=97.70$. The figure reports the true gaps, the recovered errors and each family's optimized score. The calculation is reproducible without treating the score as a percentage of people who understand the chart.

\begin{figure}[p]\centering
\includegraphics[width=\linewidth,height=.85\textheight,keepaspectratio]{reader_figures/figures/Figure_1_How_scoring_works.pdf}
\caption{\textbf{A visible group difference and an auditable score.} WDBC mean cell-nucleus area illustrates the five-gap task. \textbf{a}, Best valid score within each familiar family under the original 21-setting search. \textbf{b}, The selected raw-dot raster; each dot is one record's area summary. \textbf{c}, True malignant-minus-benign gaps and absolute recovery errors averaged over four placements. The average error is normalized before conversion to the score; rounding occurs only for display. Axes and labels lie outside the scored panels. The score measures computational recovery, not human accuracy.}\label{fig:scoring}
\end{figure}

\subsection*{Box plots succeed when the question and budget change}
The fixed-setting cross-task audit uses the same 127 comparisons and seven familiar families. Box plots belong to the winning set in 19 median comparisons, 30 interquartile-spread comparisons and two five-gap comparisons (Fig.~\ref{fig:tasks}). Ties are counted for every winning member, so column totals can exceed 127. In the original 21-setting five-gap search, box plots won none. Both the question and the opportunity to tune competing families therefore matter; the earlier zero cannot be interpreted as a verdict on box plots.

The cell-boundary example makes the distinction concrete (Fig.~\ref{fig:improvement}e). For the interquartile-spread contrast in WDBC mean concave points, a box plot scores 98.65 while mean bars score 53.74. Quartiles are explicitly encoded by the box, whereas group means do not preserve the requested spread. The implemented box uses minimum, quartiles, median and maximum; these findings do not cover every whisker convention or styling choice. The interval display also often wins median tasks because it explicitly marks the median. Popularity is neither a scoring term nor an outcome measured here.

\subsection*{Composition, ordered profiles and association broaden the candidate set}
The new extension contains 24 empirical compositions, 72 ordered heart-rate profiles and 53 paired ECG--pulse association instances derived from existing resources. These are task instances, not 149 new participants. Another 48 synthetic Dirichlet compositions stress category count and balance separately from empirical results. The extension produces 10,912 placement--observer recovery records; all decode successfully and their losses are independently recomputed from saved recovered values.

Pie charts, donuts, category bars and stacked bars compete only for shares. Under the primary observer at 256 pixels, pies win 16 of the 24 empirical cases and donuts eight. The other two families are still displayed in the complete results with zero wins. A pie represents a genuine opportunity for the composition contract, not a forced entrant into a two-group quantile question. The heart-failure example uses observed cohort label shares, not a claim about population disease risk.

For the ordered profiles, dots win 66 of 72 cases and lines six. Area and bar encodings are evaluated but do not win under this observer. For association, a 16-by-16 probability heatmap wins 44 of 53 cases and scatter plots nine. These outcomes depend on the specified recovery rules: filled profiles are read from their upper boundary, line and dot profiles from pixel centroids, and association from recovered joint pixel mass. Opaque-point overlap can alter the apparent correlation even when a person could recognize a relationship. The counts therefore describe implementations and computational observers, not universal human preferences.

Figure~\ref{fig:improvement} combines these findings into eight genuinely selected families. Different questions admit different sets, and a fair numerical comparison is made only within each panel. The starting chart is a disclosed illustrative baseline, not necessarily the strongest alternative. Full candidate scores accompany every example. Case selection is retrospective; no overall gain is estimated from the eight panels.

\begin{figure}[p]\centering
\includegraphics[width=\linewidth,height=.83\textheight,keepaspectratio]{reader_figures/figures/Figure_2_When_optimization_helps.pdf}
\caption{\textbf{Eight examples, eight computed winning families.} \textbf{a--d}, Dots, histograms, violins and heatmaps for five distributional gaps. \textbf{e}, A box plot for an interquartile-spread contrast. \textbf{f}, A pie chart for category shares. \textbf{g}, A line chart for an ordered profile. \textbf{h}, A scatter plot for paired association. The question, baseline, selected chart and scores are printed in every panel. Unlike Fig.~\ref{fig:scoring}'s full search, this audit uses one declared setting per eligible family. Scores are comparable within a panel, not across questions. Original scored pixels are enlarged; empty heatmap margins are cropped without changing colors or measurement limits. Source identities and complete scores appear in Supplementary Section S17.}\label{fig:improvement}
\end{figure}

\subsection*{A reference winner is not a guarantee at another display condition}
We freeze each case's selected family at 256 pixels, then evaluate that unchanged choice at 384 pixels with four new placements. The later-size optimum is calculated only as a hindsight comparator. Under the primary observer, a reference winner remains in the winning set in 86 of 127 median tasks, 78 spread tasks, 108 five-gap tasks, 20 of 24 empirical compositions, 66 of 72 profiles and 48 of 53 associations (Table~\ref{tab:transfer}). These descriptive counts include related features and repeated source recordings; they are not independent population replications.

\begin{table}[ht]\centering\small
\caption{\textbf{Does the frozen choice remain a winner at a new display size?} Empirical cases, primary observer. ``Retained'' means membership in the new winning set within loss tolerance $10^{-12}$. Candidate settings are not retuned. The denominator counts task instances, not independent participants.}\label{tab:transfer}
\begin{tabular}{lrrrr}\toprule
Task & Cases & Candidates & Retained & Changed\\\midrule
Median contrast & 127 & 7 & 86 & 41\\
IQR contrast & 127 & 7 & 78 & 49\\
Five gaps & 127 & 7 & 108 & 19\\
Category shares & 24 & 4 & 20 & 4\\
Ordered profile & 72 & 4 & 66 & 6\\
Correlation & 53 & 2 & 48 & 5\\
\bottomrule\end{tabular}
\end{table}

Figure~\ref{fig:limits} shows one transfer example for each of the eight selected families in Fig.~\ref{fig:improvement}. For each family and task, we select the empirical case with the largest score deficit relative to the new-size optimum; exact ties are resolved by case identifier. This deliberately emphasizes fragility and cannot estimate its prevalence. Seven illustrated choices change, whereas the illustrated pie remains best. The complete frequencies, alternative observer, failures and paired comparisons against declared fixed defaults remain available. The additional resolution changes the observed winner, not the underlying scientific question or observations.

\begin{figure}[p]\centering
\includegraphics[width=\linewidth,height=.83\textheight,keepaspectratio]{reader_figures/figures/Figure_3_When_optimization_does_not_help.pdf}
\caption{\textbf{Eight selected families face the same transfer check.} Each left-hand family was a winner at 256 pixels. Both displayed charts and scores are evaluated at 384 pixels: the frozen choice on the left and the hindsight best on the right. Seven examples reverse; the pie example retains its choice. Cases are chosen by largest transfer score deficit within each illustrated family and task, not as a representative sample. This is a display-transfer test, not external-population validation. The observer has not been validated against readers. The same chart can remain useful despite losing a close numerical rank.}\label{fig:limits}
\end{figure}

\begin{figure}[p]\centering
\includegraphics[width=\linewidth,height=.84\textheight,keepaspectratio]{reader_figures/figures/Figure_4_Question_first_chart_second.pdf}
\caption{\textbf{All familiar candidates remain visible, including non-winners.} Cells report reference winning-set memberships under the primary observer for empirical cases. A dash denotes an unimplemented task contract; zero denotes a tested family with no wins. Median, spread and five-gap columns reuse the fixed seven-family audit, so their tuning budget differs from Fig.~\ref{fig:scoring}. Ties can make a column sum exceed its case count. Mean bars are lossy diagnostic controls for marginal-distribution questions. Heatmaps use one-dimensional probability strips for the five-gap tasks and a two-dimensional probability grid for association; these are separately specified implementations. No score or winner count is pooled across questions.}\label{fig:tasks}
\end{figure}

\subsection*{Strong controls and information limits qualify every score}
The familiar-set restriction defines an application scope, not unrestricted superiority. Supplementary ECDF and direct quantile-marker controls address distributional targets explicitly. The five-marker control belongs to the winning set in 122 of 127 comparisons when the expanded 23-setting search is allowed. A 19-marker control further dominates its matched dense-quantile target. The earlier multi-task, resource-disjoint audit shows that case-specific selection does not consistently improve on a strong matched default. These results remain disclosed rather than removed to produce a more diverse set of winners.

Missing information and pixel error must also be distinguished. If two inputs have the same encoded image and calibration but different targets, no decoder can recover both exactly. For example, paired $x=(0,1,2)$ with $y=(0,1,2)$ has median within-pair change zero; pairing the same marginal $y$ values as $(1,2,0)$ changes that median to one. A marginal chart cannot distinguish those targets. The framework rejects pairing-dependent recovery from such representations instead of assigning an apparently precise score.

A small group-contrast error need not mean that both distributions are accurately represented. Common errors in the two groups can cancel. Group-level recovery losses and zero-difference controls therefore remain separate diagnostics in the distributional audit. A high score near a null contrast may add little over predicting no difference. Likewise, a small composition error under color-area counting does not establish that people estimate sectors more accurately than bar lengths. None of these limitations can be repaired by an arbitrary tie-breaking offset or an unvalidated perceptual penalty.

\section*{Discussion}
The central contribution is a reproducible question-first selection procedure whose claims can be traced to implemented targets and recovery rules. The expanded examples demonstrate that box plots and pie charts can be appropriate competitors without being forced to answer the same question. Changes in the question, candidate budget and display condition all change observed winners. An operational selector should disclose those dependencies, retain ties and distinguish unsupported tasks from tested low-scoring candidates.

The result is not a new taxonomy of visualization tasks or a claim that mathematical error alone defines communication quality. Task abstractions and human graphical-perception research supply essential foundations.\cite{brehmer2013,cleveland1984,heer2010} Our contribution is the executable connection between task declarations, chart-specific pixel recovery, scores and falsifiable transfer checks. The score can identify implementation behavior and reveal a mismatch between a question and an encoding. It cannot, without additional evidence, establish reading speed, comprehension, accessibility or scientific usefulness.

Decoder dependence is especially visible in the new experiments. Pie and donut shares can be recovered by counting calibrated colored area. This makes them strong machine-readable encodings but does not overturn human perception evidence. A calibrated probability heatmap can similarly preserve numerical mass that is visually faint.\cite{crameri2020} Point overlap, smoothing, line thickness and the chosen extraction rule can reverse a rank. Both observers use known chart geometry; neither is an independently trained or human-calibrated model. Sensitivity agreement supports only this narrow computational interpretation.

The expanded catalogue contains 13 implemented familiar families, not every widely used chart. Family names cover different constructions: a bar may encode means, category shares or ordered bin values; a heatmap may encode one- or two-dimensional probability mass. The contract, rather than the family name alone, is the reproducible unit. Geographic maps, treemaps and other contextual catalogue entries lack a tested contract here. No worldwide chart-popularity ranking is asserted.

Validation has three distinct levels. The present checks establish executable numerical consistency and rendering sensitivity. The retrospective transfer reuses source data and therefore tests display changes, not a new population. Human validation remains unperformed. Before a claim of improved scientific communication, a prospective study should freeze predictions on untouched examples, randomize task-appropriate chart assignment and measure estimation error, time and confidence. It should test selection against a fixed, strong task-matched default, including specialized distributional controls when relevant, and account for repeated participants and cases. The supplied protocol is a plan, not a completed experiment.

The benchmark remains a convenience sample of public resources. Multiple features share cohorts; BIDMC contains repeated recordings from 46 source patients, and the wrist resource contains eight participants. The 149 additional empirical instances are derived views, not additional biological replication. Source issues involving retinopathy linkage, CKD potassium values and Parkinsons identifiers remain documented. We do not infer disease mechanisms, clinical risk or treatment effects. These limits constrain the intended use to transparent computational chart evaluation pending independent validation.

\section*{Methods}
\subsection*{Sources, task instances and transformations}
The retained two-group arrays comprise WDBC 30 feature comparisons, Cleveland five, heart failure six, CKD 11, ILPD nine, retinopathy 16, HAR 40, BIDMC five and wrist five.\cite{wdbc,heart,failure,kidney,liver,retina,har,bidmc,wrist} HAR uses subject/activity medians from 30 participants; BIDMC aggregates source readings within 46 patients; wrist exercise uses early/late ECG-reference summaries within eight participants.\cite{pimentel2017,jarchi2017} Parkinsons remains excluded from the primary analysis because parsed identifier groups do not reconcile with the documented participant count. Original arrays and source-specific exclusions are unchanged.

For composition, the first canonical feature in each of nine resources defines pooled quartile boundaries; each group contributes a four-category share vector, giving 18 cases. Repeated boundaries yield empty categories rather than altered data. The first feature's valid group counts in six unpaired resources contribute six binary cohort compositions. These are observed-record proportions, not prevalence estimates. Synthetic shares use Dirichlet vectors with 2, 3, 5 or 8 categories, concentrations 0.3, 1 or 5, and four replicates per combination (48 cases; seed 20260921).

Ordered profiles use all 19 wrist exercise records and 53 BIDMC source recordings. Positive valid heart-rate values are averaged within 12 equal-duration bins. An empty bin or zero profile SD would exclude a case; no case meets those exclusions. Association uses paired positive HR and pulse readings from each BIDMC recording, systematically subsampled to at most 64 evenly spaced rows. Correlation targets refer to this displayed subset, not all original readings. No recording has undefined correlation in the final input set. Source recording counts must not be interpreted as independent patient counts.

\subsection*{Rendering, calibration and recovery}
The retained distributional renderer uses the same zero-inclusive axes, fourfold supersampling, Lanczos reduction and 0.5-pixel Gaussian blur as the earlier audit. The fixed familiar configurations are mean bars, five-number boxes and intervals, raw dots, 16-bin histograms and probability strips, and 64-grid-point violins. The fuller search permits 21 settings and is reported separately. Median and IQR targets are evaluated from saved decoded quantiles in the fixed-setting experiment; no raw target values are supplied to a decoder.

New renderings occupy a common square panel: category bars, pies, donuts with a 52\% inner-radius fraction, and stacked bars for composition; two-pixel lines, filled areas, radius-two dots and bars for profiles; radius-two opaque scatter marks and a 16-by-16 probability heatmap for association. The palette, category order and geometric settings are fixed. Selection uses 256-pixel panels at placements $0,.25,.5,.75$; transfer uses 384 pixels at $.125,.375,.625,.875$. Supersampling and blur match the retained renderer. Exact geometry is specified in the code.

For composition, the hard observer assigns non-white pixels to the closest declared category color and retains matches within RGB distance 80; normalized colored areas estimate shares. The soft observer weights assigned pixels by projected opacity. Profile readers sample the known 12 time positions: line/dot estimates use pixel centroids, and bar/area estimates use the upper recovered boundary. Profile axes include zero for all candidates; this implementation is restricted to nonnegative profiles. Association readers estimate correlation from scatter pixel mass or decoded probability-cell masses; the soft heatmap reader averages a local $3\times3$ patch. These are specified algorithms, not human models. Calibration contains geometry, axes, family and colors, never the numerical target or source observations.

\subsection*{Selection, controls and verification}
Four placement losses are averaged before score conversion. All tied minima within $10^{-12}$ are retained; a declared deterministic family order selects one representative for transfer summaries. Default comparisons use boxes for median and spread, raw dots for five gaps, category bars for shares, lines for profiles and scatter plots for association. Default and selected scores are paired only where both decode successfully; missing failures and denominators are disclosed. These defaults are retrospective fixed controls, not independently established best practice.

New fixture checks cover all ten task--family contracts under both observers, blank-image failure and profile unit-rescaling invariance. Every one of the 10,912 new losses is recomputed from saved recovered outputs and independently calculated targets; the maximum absolute discrepancy is $4.45\times10^{-16}$. No new trial fails. Existing quantile-decoder checks, original reference parity, source sensitivities and specialized-control audits remain in the Supplement. The transfer selection never uses later-size losses; the later oracle is descriptive only.

The extension is retrospective and not preregistered. Its new recipes and settings were written before the new rendering run, while earlier results and source data had already been inspected. Both empirical and synthetic outputs are fully retained, with no winner quota or outcome-based data correction. Figure 2 examples were selected after inspection to demonstrate eight actual winning families. Figure 3 deliberately selects the largest empirical transfer deficit within those family--task pairs. Protocol, inputs, code, per-placement recovered values and checksums accompany the release.

\section*{Data availability}
Sources are available through the cited UCI and PhysioNet records. Prepared arrays, supplied source copies, transformations, acquisition records and checksums accompany the reproducibility archive, subject to original terms documented in \texttt{DATA\_PROVENANCE.md}. New task inputs and synthetic labels are in \texttt{task\_first/results/inputs.json}. No new participants or human-reader responses were collected.

\section*{Code availability}
The archive contains pinned dependencies and runnable examples. \texttt{score\_task.py} requires a declared question and evaluates its implemented candidates; \texttt{score\_my\_data.py} retains the two-group workflow. \texttt{reproduce\_task\_first.py} rebuilds the extension, and the earlier audits remain executable. Figure sources, all candidate scores and enlarged examples are included. Repository: \url{https://github.com/Elgendi/two-sample-chart-selection}.

\section*{Author contributions}
M.E. conceived the study, developed the methodology, performed the analyses, interpreted the results, and wrote and revised the manuscript.
\section*{Competing interests}
M.E. serves on the editorial boards of Communications Medicine and npj Biosensing. The author declares no other competing interests.
\begingroup\scriptsize\setlength{\bibsep}{0pt}\bibliographystyle{viznumbered}\bibliography{references}\endgroup
\end{document}
